% ~215 words. The claim ladder in the order the paper makes it. Must match the registered
% abstract, so change this and the contribution list of p01_intro together.
Genomic language models are released pretrained and fine-tuned locally, increasingly on cohort data.
We ask what happens to the individuals in that cohort, and measure it on six released models rather
than on a backbone of our own.

Every released masked checkpoint we tested reproduces planted records verbatim at a held-out loss a
practitioner would accept: the cheapest costs seven thousandths of a bit and returns 30\% of records.
On real chromosome-22 haplotypes, each carried by at most five of 2{,}548 individuals, a record
present as a single copy comes back letter for letter about one time in ten.

Reproduction is the wrong alarm. At checkpoints where verbatim recovery reads exactly zero in every
seed, the same model ranks one individual's haplotype first against a hundred real rare segments
0.92 to 0.96 of the time, with none of 350 never-trained controls ever identified. Downstream task
score moves one to two per cent across the same range.

Whether a memorised record can be reached at all is set by how many token slots it occupies, not by
how much the model knows: models differing only in tokenisation answer any single hidden nucleotide
identically and yet return 92\% or 0\% of whole records. No single extraction test is correct across
tokenizers, which is why this has been missed. Code, measurements and logs are public.
